\scanpage{55}
% @ lines 1–2 / text
On trouve, en notant que tt ne dépend que de la section correspondante de $G=\check{\mathbb V}(M^*)/G_m$ :
% @ lines 2–3 / text
\uline{Prop} Pour que $g=\bar u\in G(S)$ invarie $Y$, il f. et s. qu'il invarie $A$ i.e. que \deleted{\unsure{$\operatorname{int}(u)A=A$}.}
% @ line 4 / display
\[
 \operatorname{int}(u)A\ \bigl(\overset{\mathrm{dfn}}{=}uAu^{-1}\bigr)=A
\]
% @ lines 5–8 / text
Quand il en est ainsi, l'opération de $g$ sur $Y$ \deleted{est} s'identifie via (23) à l'op. de $g$ \inserted{(ou de $u$)} sur $\operatorname{Spec}(A)$ -- donc elle est triviale ssi $g$ opère trivialement sur $A$ i.e ssi $u\in\Gamma(S,A)$.

% @ lines 9–11 / text
Bien entendu, ceci caractérise $A$ \deleted{\unsure{via|comme}} (engendré comme sous-fibré vectoriel par les sections locales de $A^*$ ...).

% @ line 12 / text
À \deleted{\unsure{tout}} $A$, associons le ss-grpe
% @ lines 13–14 / display
\[
 \underset{\substack{\shortparallel\\\check{\mathbb V}(A)^*/G_m}}{\unsure{\mathscr J|\mathscr T|\mathcal T}}\subset G\tag*{(25)}
\]
% @ lines 15–20 / text
c'est un sous-grpe \uline{commutatif} \uline{lisse} \uline{connexe} de dim \uline{relative} $1$. Bien entendu, la correspondance $A\mapsto\mathscr J$ est injective (\deleted{\unsure{$\mathbb V$}} $A$ étant engendré par \deleted{\unsure{l'}} les sections de $\check{\mathbb V}(A)^*$, qui est l'image inverse de $\mathscr J$ par l'hom can $\check{\mathbb V}(M)^*\longrightarrow G$.
% @ lines 21–24, continues / text
D'autre part, \inserted{si $2$ est inv. sur $S$,} tout $\mathscr J$ ayant les propriétés énumérées (sans $=$ la \unsure{commutati...|commutativité...|connexité...}) \deleted{-- appelons les les ss-grpes \uline{réguliers} de $G$ -- proviennent de tels $A$} -- i.e. de ss. fibrés inv
% archive sheet number: 54
